\chapter{Wyniki i analiza}

\section{Wyniki scenariusza bazowego}
W scenariuszu bazowym (mirrorowanie 10\% ruchu) wersja SMF v2 przetwarzała ruch bez naruszenia obsługi produkcyjnej. Zaobserwowano:
\begin{itemize}
    \item stabilną średnią latencję poniżej progu SLO,
    \item brak błędów HTTP 4xx/5xx dla monitorowanego okna pomiarowego,
    \item poprawną separację danych pomiędzy \texttt{udr-prod} i \texttt{udr-shadow}.
\end{itemize}

Iter8 zwracał werdykt \texttt{pass}, co oznacza spełnienie kryteriów jakości dla wersji kandydackiej.

\section{Wyniki scenariusza przeciążeniowego}
Przy zwiększonym obciążeniu zaobserwowano wzrost czasu odpowiedzi SMF v2, jednak bez przekroczenia progu 100 ms w oknie referencyjnym. Wariant ten potwierdził, że:
\begin{itemize}
    \item polityka mirrorowania nie destabilizuje ścieżki produkcyjnej,
    \item komponent shadow wymaga odrębnego planowania zasobów,
    \item utrzymanie zapasu CPU i pamięci jest kluczowe dla wiarygodnej walidacji.
\end{itemize}

\section{Wyniki scenariusza błędów}
W scenariuszu kontrolowanym, gdzie wprowadzono sztuczny błąd w logice SMF v2, odsetek błędów HTTP wzrósł powyżej przyjętego limitu. Iter8 zwrócił werdykt \texttt{fail}, co skutecznie zablokowało promocję wersji.

Wynik ten potwierdza praktyczną użyteczność automatycznych \emph{quality gates} i zasadność stosowania metryk pochodzących bezpośrednio z Service Mesh, co jest spójne z praktykami nowoczesnych strategii canary/shadow~\cite{fink2017canary}.

\section{Analiza korzyści i kosztów}
Najważniejsze korzyści wdrożenia Shadow Deploy:
\begin{itemize}
    \item redukcja ryzyka wdrożeń na krytycznych funkcjach 5G Core,
    \item szybsze wykrywanie regresji wydajnościowych i funkcjonalnych,
    \item większa przewidywalność procesu release management.
\end{itemize}

Najważniejsze koszty i wyzwania:
\begin{itemize}
    \item dodatkowe zużycie zasobów przez ścieżkę shadow,
    \item większa złożoność konfiguracji Service Mesh i obserwowalności,
    \item konieczność utrzymania spójnych danych testowych w \texttt{udr-shadow}.
\end{itemize}

\section{Dyskusja}
Uzyskane rezultaty wskazują, że architektura Shadow Deploy jest szczególnie użyteczna dla funkcji kontrolnych 5G Core, gdzie koszt błędu produkcyjnego jest wysoki. W praktyce operatora podejście to może stanowić etap pośredni między testami laboratoryjnymi a pełnym wdrożeniem produkcyjnym.

Jednocześnie skuteczność metody zależy od jakości zdefiniowanych metryk i progów SLO. Zbyt liberalne progi zwiększają ryzyko przepuszczenia regresji, natomiast zbyt restrykcyjne progi mogą blokować poprawne wdrożenia. W praktyce użyteczne jest łączenie oceny metrycznej z analizą odpowiedzi przy pomocy narzędzi takich jak Diffy~\cite{diffyGithub}.

\section{Automatycznie wygenerowane metryki}
\label{sec:auto-metrics}

\IfFileExists{generated/metrics/scenario.tex}{
    \input{generated/metrics/scenario.tex}
    Dane pomiarowe pochodza z automatycznego eksportu metryk dla scenariusza \texttt{\MetricsScenario} (\texttt{\MetricsTimestamp}), w oknie \texttt{\MetricsWindowStart} -- \texttt{\MetricsWindowEnd} przy kroku zapytan \texttt{\MetricsStep}. Eksport zawiera \MetricsExportsCount{} zapytan, a liczba brakujacych wynikow to \MetricsMissingCount{}.
}{
    Brak wygenerowanych metryk. Uruchom skrypt \texttt{archive/tools/build-with-metrics.sh --metrics-dir <sciezka-do-exportu>}, aby dolaczyc dane i wykresy.
}

\ifmetricsplots
    \IfFileExists{generated/metrics/request_rate.tsv}{
        \begin{figure}[htbp]
            \centering
            \begin{tikzpicture}
                \begin{axis}[
                    width=\textwidth,
                    height=0.27\textheight,
                    date coordinates in=x,
                    xticklabel style={rotate=35,anchor=east,font=\footnotesize},
                    x tick label style={/pgf/number format/1000 sep=},
                    xlabel=Czas (UTC),
                    ylabel={Request rate [rps]},
                    grid=both
                ]
                    \addplot[
                        thick,
                        blue
                    ] table [col sep=tab, x=ts_iso, y=value_rps] {generated/metrics/request_rate.tsv};
                \end{axis}
            \end{tikzpicture}
            \caption{Request rate w czasie.}
        \end{figure}
    }{}

    \IfFileExists{generated/metrics/latency_p95_p99.tsv}{
        \begin{figure}[htbp]
            \centering
            \begin{tikzpicture}
                \begin{axis}[
                    width=\textwidth,
                    height=0.27\textheight,
                    date coordinates in=x,
                    xticklabel style={rotate=35,anchor=east,font=\footnotesize},
                    x tick label style={/pgf/number format/1000 sep=},
                    xlabel=Czas (UTC),
                    ylabel={Opoznienie [s]},
                    legend pos=north west,
                    grid=both
                ]
                    \addplot[
                        thick,
                        red
                    ] table [col sep=tab, x=ts_iso, y=p95_seconds] {generated/metrics/latency_p95_p99.tsv};
                    \addlegendentry{p95}
                    \addplot[
                        thick,
                        magenta
                    ] table [col sep=tab, x=ts_iso, y=p99_seconds] {generated/metrics/latency_p95_p99.tsv};
                    \addlegendentry{p99}
                \end{axis}
            \end{tikzpicture}
            \caption{Latencja p95/p99 w czasie.}
        \end{figure}
    }{}

    \IfFileExists{generated/metrics/error_rate.tsv}{
        \begin{figure}[htbp]
            \centering
            \begin{tikzpicture}
                \begin{axis}[
                    width=\textwidth,
                    height=0.27\textheight,
                    date coordinates in=x,
                    xticklabel style={rotate=35,anchor=east,font=\footnotesize},
                    x tick label style={/pgf/number format/1000 sep=},
                    xlabel=Czas (UTC),
                    ylabel={Error rate [\%]},
                    grid=both
                ]
                    \addplot[
                        thick,
                        magenta
                    ] table [col sep=tab, x=ts_iso, y=value_percent] {generated/metrics/error_rate.tsv};
                \end{axis}
            \end{tikzpicture}
            \caption{Odsetek bledow 5xx w czasie.}
        \end{figure}
    }{}

    \IfFileExists{generated/metrics/prod_shadow_counts.tsv}{
        \begin{figure}[htbp]
            \centering
            \begin{tikzpicture}
                \begin{axis}[
                    width=\textwidth,
                    height=0.27\textheight,
                    date coordinates in=x,
                    xticklabel style={rotate=35,anchor=east,font=\footnotesize},
                    x tick label style={/pgf/number format/1000 sep=},
                    xlabel=Czas (UTC),
                    ylabel={Liczba requestow},
                    legend pos=north west,
                    grid=both
                ]
                    \addplot[
                        thick,
                        green
                    ] table [col sep=tab, x=ts_iso, y=prod_requests] {generated/metrics/prod_shadow_counts.tsv};
                    \addlegendentry{prod}
                    \addplot[
                        thick,
                        cyan
                    ] table [col sep=tab, x=ts_iso, y=shadow_requests] {generated/metrics/prod_shadow_counts.tsv};
                    \addlegendentry{shadow}
                \end{axis}
            \end{tikzpicture}
            \caption{Porownanie liczby requestow prod vs shadow.}
        \end{figure}
    }{}

    \IfFileExists{generated/metrics/pod_cpu_usage.tsv}{
        \begin{figure}[htbp]
            \centering
            \begin{tikzpicture}
                \begin{axis}[
                    width=\textwidth,
                    height=0.27\textheight,
                    date coordinates in=x,
                    xticklabel style={rotate=35,anchor=east,font=\footnotesize},
                    x tick label style={/pgf/number format/1000 sep=},
                    xlabel=Czas (UTC),
                    ylabel={CPU [cores]},
                    grid=both
                ]
                    \addplot[
                        thick,
                        blue
                    ] table [col sep=tab, x=ts_iso, y=cpu_cores] {generated/metrics/pod_cpu_usage.tsv};
                \end{axis}
            \end{tikzpicture}
            \caption{Trend zuzycia CPU dla kluczowych podow (suma).}
        \end{figure}
    }{}

    \IfFileExists{generated/metrics/pod_memory_usage.tsv}{
        \begin{figure}[htbp]
            \centering
            \begin{tikzpicture}
                \begin{axis}[
                    width=\textwidth,
                    height=0.27\textheight,
                    date coordinates in=x,
                    xticklabel style={rotate=35,anchor=east,font=\footnotesize},
                    x tick label style={/pgf/number format/1000 sep=},
                    xlabel=Czas (UTC),
                    ylabel={Pamiec [MiB]},
                    grid=both
                ]
                    \addplot[
                        thick,
                        black
                    ] table [col sep=tab, x=ts_iso, y=memory_mib] {generated/metrics/pod_memory_usage.tsv};
                \end{axis}
            \end{tikzpicture}
            \caption{Trend zuzycia pamieci dla kluczowych podow (suma).}
        \end{figure}
    }{}
\else
    \paragraph{Uwaga}
    Pakiet \texttt{pgfplots} nie jest dostepny, dlatego wykresy metryk nie zostaly wyrenderowane.
\fi

